\section{方法}

\subsection{信息型动作词与攻击性关系的切分}

记 $A$ 为「政治攻讦」关系集合（$|A|=193$），$F$ 为「敌对冲突」（$|F|=321$），
$R$ 为「君臣从属」（$|R|=397$）。令 $\mathcal\{V_{act}\}=\{$告发、告密、弹劾、诬告、监视、构陷、谗、上书$\}$
为信息型动作词表，约定 $\mathrm{I_{act}}(r)\iff\exists v\in\mathcal{V_{act}}:\ v$ 出现在 $r$ 的 $\texttt{rel}$ 字段中。
对任意关系子集 $S$，定义信息型命中率 $\eta(S)=100\,|\{r\in S:\mathrm\{I_{act}\}(r)\}|/|S|$。
本文的第一个对比量就是 $\eta(A)$、$\eta(F)$、$\eta(R)$ 三者的差。

\subsection{时间代理：以书的七部为刻度}

$\texttt{relations}[].chapter$ 存的是章节号而非年份，且 40.4\% 的关系无章节归属。本文因此**不用年份**，
改用两个更稳的刻度：

\textbf{(a) 分部占比。}记 $p\in\{1,\dots,7\}$ 为书的分部，$\Dcal_p$ 为第 $p$ 部中带章节归属的全部关系
（分母），$A_p$ 为该部中「政治攻讦」的条数（分子），定义
$\pi_p(A)=100\,A_p/\Dcal_p$。分部是书的物理顺序，不受抽取年份影响。

\textbf{(b) 三年段代理（仅作旁证）。}先把章节映射到该章事件年份的中位数，再按
1368--1450 / 1451--1550 / 1551--1644 三段归并。这一路覆盖 170/193 条攻讦关系（其余 23 条缺年份），
只能作方向性旁证，不能替代 (a)，正文结论一律以 (a) 为准。

两个口径不一致时以 (a) 为准，这是本文的先验约定，为了可复现而写下，不是为了保结论而挑选。

\subsection{机构与代理人的辨认}

$\Ccal$ 为全部人物。定义
$\Igc=\{c\in\Ccal : \exists v\!\in\!\{$锦衣卫,东厂,西厂,内行厂$\}$,\ v 出现在 $c$ 的 $\texttt{role}$ 或
$\texttt{faction}$ 中$\}$（厂卫系，$|\Igc|=29$），
$\Iec=\{c\in\Ccal : \exists v\!\in\!\{$太监,宦官$\}$,\ v 出现在 $c$ 的 $\texttt{role}$ 或
$\texttt{faction}$ 中$\}$（内廷宦官，$|\Iec|=51$）。二者**不互相排斥**（宦官可兼有厂卫身份），
本文从不把两者相减或相加成一个「合并机构」数。

对每个 $c$ 计算无向度数 $\Deg(c)$（$\texttt{relations}[]$ 的 $from$、$to$ 各计一次）与出场章节数 $\Chap(c)$，
得到全库均值 $\overline{\Deg}(\Ccal)=3.78$、$\overline{\Chap}(\Ccal)=2.22$，
以及相对比值 $\mathrm{ratio}(S)=\overline{\Deg}(S)/\overline{\Deg}(\Ccal)$。

\subsection{稳健性设计}

三重：
\begin{enumerate}
  \item \textbf{分母敏感性}：所有分部占比同时用「该部全部关系」和「该部带章节归属的关系」两种分母算一遍，
        结论取两者交集（正文结果在两者下同号）。
  \item \textbf{派系口径敏感性}：同阵营／跨阵营判定先用完整派系串（如「明朝·阉党」与「明朝·文官」判为跨阵营），
        另用折叠到根类（都归「明朝」）的口径复算一遍。折叠口径会把 76.7\% 的攻讦判成「同阵营」，
        因为超过九成人物挂在同一根类下——这个数字不可用，正文只报完整串口径。
  \item \textbf{端点集中度}：用基尼系数与 Top-10 占比刻画攻讦端点分布，防止结论只由个别人物撑起。
\end{enumerate}

\textbf{可复现性}：全部数字由 `papers\textunderscore{}pilot/info/\_stats\_py` 写入
`.dump/\_stats\_info\.json`，三张图各有一个独立脚本，数据一律从该 JSON 读取，脚本内不写死数字。

\textbf{首要效度威胁（先声明）。}关系类别与派系标签都由同一份语料反查导出，
「攻讦随叙事上升」这一现象不能被排除是**抽取顺序的副产品**：作者写到后面章节时对人物关系更熟，
自然抽出更多攻讦边。这是本篇的首要限制，讨论节会回到这一点。
